\section{Exceptional Scaler, Conformal Theory, Strings and Mass Classification}
\Needspace{7\baselineskip}
\subsection{Common preimage and coordination of the \(2\) readings of the exceptional corridor}
The relation between masses and exceptional symmetries does not begin in an action of the Monster on a table of particles. It begins in the same enriched state of the TPK. Let

\[
x\in\mathcal X_{\rm TPK}^{\rm enr}.
\]
The Archimedean reading contracts compatible cylinders and preserves the real value with its genealogy:

\[
\mathcal P_{\rm arq}(x)\in\mathbb R.
\]
The exceptional reading preserves boundary incidences and produces a class in the projective limit:

\[
\mathcal P_{\rm exc}(x)=Q_\infty(x).
\]
Naturality requires, for every refinement \(n\ge m\),

\[
r_{n,m}\circ Q_n=Q_m\circ p_{n,m}.
\]
The joint application

\[
\mathcal P_{\rm joint}:x\longmapsto
\bigl(\mathcal P_{\rm arq}(x),\mathcal P_{\rm exc}(x)\bigr)
\]
is the positive form of the double projection. It does not deny the Monster--mass connection: it types it. The mass branch and the exceptional branch are images of the same preimage with common orientation, grading, carry, memory, and holonomy. Nor does it assert that the Monster acts literally on decimal digits or that a Moonshine coefficient is by itself a mass.

\Needspace{7\baselineskip}
\subsection{Construction of the exceptional corridor}
The ternary incidence prolongs the state by the chain

\[
C_3
\longrightarrow A_2^{12}
\longrightarrow N(A_2^{12})
\longrightarrow \Lambda_{24}
\longrightarrow V^\natural
\longrightarrow \mathbb M,
\]
where \(\Lambda_{24}\) is the Leech lattice, \(V^\natural\) the Moonshine vertex algebra, and \(\mathbb M\) the Monster group. The step does not consist in juxtaposing classical names: each arrow must preserve the incidence, grading, and refinement compatibility inherited from the state.

The ternary trace is written

\[
t_3(q)=q^{-1}
\prod_{n\ge1}(1+q^n+q^{2n})^{-12}.
\]
The first nontrivial defect satisfies

\[
D_{\rm APP}=[q]t_3=v_\alpha^2=54,
\]
and the first coefficient of the Monster's character preserves the decomposition

\[
196884=54+10\cdot3^9.
\]
These identities transport gradation and multiplicity. They do not assign a mass in MeV. Their function in the theory of masses is to classify the space on which the operator acts before scalarizing it.

\Needspace{7\baselineskip}
\subsection{Fiberwise matching with the mass law}
The material branch of the same state produces

\[
x\longmapsto
\bigl(F_X,\sigma_X,\eta_X,\widehat M_X\bigr),
\]
where \(F_X\) is a point, orbit, fiber, or multisection; \(\sigma_X\in\mathbb Z^5\) is the signature; \(\eta_X\in\mathcal E_{90,120}\) is the refinement; and \(\widehat M_X\) is the mass operator. The joint reading is then

\[
\mathfrak J_{\rm mass,exc}(x)
=\bigl(F_X,\sigma_X,\eta_X,\widehat M_X,Q_\infty(x)\bigr).
\]
Two states with the same scalar mass value may remain distinct because of their exceptional component; two states with the same exceptional character may have different mass operators. The correct classification is the matching of fibers, not the isolated equality of numbers.

If an exceptional representation \(G\) acts in a physical fiber \(\mathcal H_X\), compatibility with mass requires

\[
U_g\widehat M_XU_g^{-1}=\widehat M_X,
\qquad g\in G.
\]
The decomposition

\[
\mathcal H_X=\bigoplus_{\lambda}
\mathcal H_{X,\lambda}
\]
separates exceptional multiplets. On an irreducible summand, commutation forces the compression of \(\widehat M_X\) to be scalar. This is a legitimate route to scalarization by symmetry. The scalar value continues to arise from the character, the towers, the dimensional section, and the coupling; the dimension of the representation determines only the degeneracy.

\Needspace{7\baselineskip}
\subsection{Fock and vertex algebra}
The prolongation of a route to a many-particle space has two principal completions. The exterior completion produces CAR, Pauli exclusion, and fermionic statistics; the symmetric completion produces CCR and bosonic statistics. A vertex algebra adds an operator product, vacuum, translation, and grading:

\[
V=\bigoplus_{n\ge0}V_n,
\qquad
L_0|_{V_n}=nI.
\]
The character

\[
\operatorname{ch}_V(q)
=\operatorname{Tr}_V q^{L_0-c/24}
\]
counts states by degree. A multiplicity in \(V_n\) is not a mass. To
convert the grading into energy, a Hamiltonian and a length or time section
are required. In a conformal realization on a cylinder of circumference
\(L\),

\[
H_{\rm CFT}
=\frac{2\pi\hbar_{\rm int}v}{L}
\left(L_0+\bar L_0-\frac{c}{12}\right).
\]
The physical mass or gap appears only after fixing \(v\), \(L\), the boundary conditions, and the map from the TPK route to the conformal modules. The role of the exceptional corridor is to organize the degeneracies and composition rules of those modules.

\Needspace{7\baselineskip}
\subsection{String spectrum as a posterior dimensional realization}
Let \(\alpha'\in\mathcal L_L^2\) be the squared length representing inverse tension in the declared naturalized string chart. In a circular compactification of radius \(R\), the left- and right-moving momenta are

\[
p_L=\frac{m}{R}+\frac{wR}{\alpha'},
\qquad
p_R=\frac{m}{R}-\frac{wR}{\alpha'},
\]
with \(m,w\in\mathbb Z\). In the natural units of the realization, the closed spectrum satisfies

\[
M^2=p_L^2+\frac{4}{\alpha'}(N_L-a_L)
=p_R^2+\frac{4}{\alpha'}(N_R-a_R),
\]
together with the level condition

\[
N_L-N_R=a_L-a_R-mw.
\]
The involution

\[
(m,w,R)\longleftrightarrow
\left(w,m,\frac{\alpha'}{R}\right)
\]
preserves the spectrum: it is the realization of T-duality. HMT first supplies the bidirectional orientation, memory, and sheet pairing; string theory adds tension, radius, oscillators, level, and boundary conditions. A string mass is not obtained by replacing \(N_L,N_R\) with heights of the tower \(\langle90,120\rangle\). Any identification between the two gradings requires an intertwiner that preserves sum, orientation, and character.

\Needspace{7\baselineskip}
\subsection{M-Theory Lattice Screens: Grading, Compactification, and Transport of the Residual Memory}
The internal screens

\[
12\longrightarrow11\longrightarrow10,
\qquad
26\longrightarrow24,
\]
preserve the dodecaphase descent, the elimination of a redundant direction, and the transverse screen. In a realization of M-theory, the fields

\[
g_{MN},\qquad C_{MNP},\qquad\Psi_M
\]
must be obtained through an intertwiner that preserves gauge, supersymmetry, constraints, and inner product. Numerical displays do not replace that map.

The classification of masses is transported to this branch by means of

\[
\text{route}
\longrightarrow\text{genealogical record}
\longrightarrow\text{signature}
\longrightarrow\text{character}
\longrightarrow\text{operator}
\longrightarrow\text{compactified spectrum}.
\]
Kaluza--Klein, winding, membrane, and oscillatory-excitation modes are distinct states even when they share a mass. The graviton of a closed-string realization, if it appears, is a massless collective mode of the geometric sector; it does not thereby become an elementary seed of the TPK or invalidate the formulation of gravity as holonomy and torsion.

\Needspace{7\baselineskip}
\subsection{Ordinary and Hypothetical Realizations Induced by the Exceptional Runner}
The exceptional branch allows distinguishing four sources of multiplicity:

\begin{enumerate}
\item cellular or route multiplicity in the atlas;
\item degeneracy of an exceptional representation;
\item Fock occupation;
\item conformal, oscillatory, or compactified level.
\end{enumerate}
Confusing them produces false companions, false generations, or duplicated masses. A twin, supersymmetric, dark, or exotic state constitutes a new species only when its quintuple

\[
(F_X,Q_\infty(x),\rho_X,\widehat M_X,P_X)
\]
differs physically, and the Hamiltonian has an additional pole or persistent section. A coincidence of dimension, character, or mass does not suffice.

Conversely, a degeneracy protected by the exceptional corridor is a structural prediction: perturbations preserving the symmetry cannot split the multiplet. An observed splitting requires identifying the symmetry breaking, coupling, or chart that produces it.

\Needspace{7\baselineskip}
\subsection{Positivity, closure, and stability conditions of topological realizations}
The connection between the exceptional corridor and masses is falsified in a concrete realization if any of the following occurs:

\begin{itemize}
\item naturality \(r_{n,m}Q_n=Q_mp_{n,m}\) fails;
\item the exceptional label changes when the external masses are permuted;
\item a Moonshine multiplicity is used as a mass value without a Hamiltonian;
\item the mass operator does not commute with the symmetry asserted to be preserved;
\item a tower height is identified with a string level without an intertwiner;
\item the level condition or T-duality invariance fails;
\item a dimensional screen is declared M-theory without constructing its fields;
\item the existence of a spin-two mode is used to redefine the preceding
internal gravity.

\end{itemize}
These controls keep the connection positive: there is a common state and a joint classification; what is prohibited is replacing the arrows with a numerical analogy.

\Needspace{7\baselineskip}
\subsection{Mathematical construction of material classes from routes, signatures, and representations}
The double realization, projective naturality, the chain \(A_2^{12}\)--Leech--\(V^\natural\)--Monster, the trace \(t_3\), the identities \(D_{\rm APP}=54\) and \(196884=54+10\cdot3^9\), and the dimensional screens belong to the APP--TRIT--TPK mathematical architecture developed previously. The formulation of \(\mathfrak J_{\rm mass,exc}\), the separation of the four multiplicities, and the foregoing controls make its function within an autonomous theory of masses explicit. The conformal and string formulas are subsequent dimensional realizations with their parameters declared; they are not used retroactively to generate APP, TRIT, TPK, the character, or the towers.
